\section{Struktura databáze}
\label{sec:database}

při návrhu databáze bylo klíčové rozhodnout jak bude databáze vypadat, tak aby efektivně pojmula velké množství dat a zároven aby umožýnovala rychlé provádění analytických dotazů.

\subsection{Realční vs dokumentová}
\label{subsec:database:relational_vs_document}
první věc o které muselo být rozhodnuto je zda vybrat dokumentovnou nebo relační databázi. Ikdyž dokumentová databáze nabízí vetší flexibilitu dat a umožnuje ukládat celé json odpovědi bez striktní struktury bylo rozhodnuto pro relační databázi

z analytické povahy dotazů a potřeby agregovat data přes více entit. Relační databáze umožnuje efektivní spojování tabulek, seskupování a další analytické operace, které jsou optimalizovány v relačních databáízích

taktéž relační databáze zajištuje díky primárním a cizím klíčům integritu dat a je předem definovaná struktura což umožnuje validaci a konzistenci dat a to zjednodušuje logiku na straně aplikace.

% odděleníá labels a events

\subsection{Rozdělení labelů a eventů}
labely a eventy jsou rozděleny z důvodu málo protínajících se atributů na to aby byly vyřešeny podobně jako u tabulky issues kde jediná změna mezi issue a pr je příznak toho zda se jedná o pr nebo issue. labely a eventy mají řadu odlišných atributů a proto je lepší je rozdělit do dvou tabulek.

\subsection{Diagram databáze}

schéma má tři hlavní logické části. první je struktura kontributorů a do jakého týmu patří pokud patří, druhá část je aktuální stav pr a issues a třetí část je historie změn od issues a to jaké eventy se staly a jaké labely byly přidány nebo odebrány a taktéž jaké soubory se změnily v daném pr pokud už je pr merged nebo closed.

centrem schématu je tabulka issues která obsahuje aktuální stav issue a pr. tato tabulka je společná jak pro issue tak pro pr z důvodu že i na straně githubu je každý pr i issue a pravděpodobně se rozděluje pouze za pomocí příznaku. například u endpointu pro získání issues se rozdíl mezi pr a issue pozná primárně pomocí přítomnosti pole \code{pull\_request} v objektu entity v odpovědi. kvůli této jediné změně je implementován podmíněný index na této tabulce. konkrétně dva jeden kde pr = true a druhý kde ispr = false. tento index effektivně rozděluje hledání v tabulce na dvvě "poloviny"a tím zmenšuje index a urychluje vyhledávání.

do historických entit se pouze přidávají události a nikdy se nemění. ostatní záznamy v tabulkách jsou aktualizovány podle posledního známého stavu.

vztah mezi týmem a kontributorem je realizován jako m:n přes spojovací tabulku contributors\_teams protože každý člen může být ve více týmech a samozřejmě tým má více členů.

ohledně souborů je logováno pouze kdo upravil daný soubor neloguje se jaké změny byly provedeny. to je dostačující protože stačí nám na jaké části projektu se daný kontributor podílel a nejde nám o konkrétní funkcionalitu kterou upravil.
